\chapter{Clase 20}

\section{Aplicaciones}

\subsection{Ondas en interfaces elásticas}

Cuando una onda elástica encuentra una interfaz entre dos materiales, parte de su energía puede reflejarse y parte transmitirse. Como los medios pueden tener propiedades distintas, durante el proceso también puede ocurrir conversión entre ondas longitudinales (P) y transversales (S).

\subsubsection{Ondas planas}

Representamos cada onda plana mediante
\[
\vec u_\alpha(\vec r,t)
=\vec U_\alpha e^{i(\vec k_\alpha\cdot\vec r-\omega t)},
\]
donde $\vec U_\alpha$ es su amplitud vectorial, $\vec k_\alpha$ su vector de onda y $\omega$ su frecuencia angular. El subíndice $\alpha$ identifica la onda incidente, reflejada o transmitida; en un sólido, cada una puede ser longitudinal (P) o transversal (S).

\subsubsection{Condiciones en la interfaz}

Consideremos una interfaz plana y estacionaria en $z=0$ entre los medios 1 y 2, perfectamente adheridos entre sí. La continuidad del desplazamiento exige que el desplazamiento total en el medio 1 coincida con el del medio 2:
\[
\vec u_{\mathrm{inc}}^{(1)}+\vec u_{\mathrm{ref}}^{(1)}
=\vec u_{\mathrm{trans}}^{(2)}
\qquad (z=0).
\]
Además, debe ser continua la tracción sobre la interfaz:
\[
\left(\boldsymbol{\sigma}_{\mathrm{inc}}^{(1)}
+\boldsymbol{\sigma}_{\mathrm{ref}}^{(1)}\right)\cdot\hat e_z
=\boldsymbol{\sigma}_{\mathrm{trans}}^{(2)}\cdot\hat e_z
\qquad (z=0).
\]
Aquí cada término puede representar la suma de los modos P y S presentes en ese medio. Estas condiciones, junto con la coincidencia de fase para todo $x$, $y$ y $t$, determinan los ángulos y amplitudes de las ondas reflejadas y transmitidas.

En particular, la frecuencia y las componentes del vector de onda paralelas a la interfaz se conservan:
\[
\omega_{\mathrm{inc}}=\omega_{\mathrm{ref}}=\omega_{\mathrm{trans}},
\qquad
k_{x,\mathrm{inc}}=k_{x,\mathrm{ref}}=k_{x,\mathrm{trans}},
\qquad
k_{y,\mathrm{inc}}=k_{y,\mathrm{ref}}=k_{y,\mathrm{trans}}.
\]
Esta es la forma ondulatoria de la ley de Snell. Las componentes normales a la interfaz pueden cambiar, y la conversión de modo permite, por ejemplo, que una onda P incidente genere ondas P y S reflejadas o transmitidas.

\begin{figure}[H]
\centering
\begin{tikzpicture}[
  >=Latex,
  x=1cm,
  y=1cm,
  line cap=round,
  line join=round,
  font=\small,
  incidente/.style={->, blue!75!black, line width=1.5pt},
  reflejada/.style={->, teal!70!black, line width=1.5pt},
  transmitida/.style={->, orange!85!black, line width=1.5pt},
  frentei/.style={draw=blue!45, line width=0.9pt},
  frenter/.style={draw=teal!45, line width=0.9pt},
  frentet/.style={draw=orange!55, line width=0.9pt}
]

\fill[blue!4] (-5.1,0.08) rectangle (5.1,3.25);
\fill[orange!5] (-5.1,-2.75) rectangle (5.1,-0.08);
\draw[black!75, line width=1.3pt] (-5.1,0) -- (5.1,0);
\node[anchor=north east, fill=white, inner sep=2pt] at (-4.75,3.02)
  {medio 1};
\node[anchor=south east, fill=white, inner sep=2pt] at (-4.75,-2.52)
  {medio 2};
\node[anchor=west, fill=white, inner sep=2pt] at (3.0,0.12)
  {interfaz $z=0$};

\draw[densely dashed, gray!75] (0,-2.35) -- (0,2.9);
\draw[->, gray!75, line width=1pt] (0.12,1.85) -- (0.12,2.65)
  node[above] {$+z$};
\draw[->, gray!75, line width=1pt] (-4.3,-2.25) -- (-3.2,-2.25)
  node[right] {$x$};

\draw[frentei] (-2.42,2.31) -- (-1.98,2.83);
\draw[frentei] (-1.77,1.54) -- (-1.33,2.06);
\draw[frentei] (-1.12,0.78) -- (-0.68,1.30);
\draw[incidente] (-2.08,2.43) -- (-0.25,0.30);

\draw[frenter] (0.33,0.90) -- (0.77,0.38);
\draw[frenter] (0.93,1.60) -- (1.37,1.08);
\draw[frenter] (1.53,2.30) -- (1.97,1.78);
\draw[reflejada] (0.25,0.30) -- (2.08,2.43);

\draw[frentet] (0.35,-0.78) -- (0.75,-0.32);
\draw[frentet] (0.95,-1.28) -- (1.35,-0.82);
\draw[frentet] (1.60,-1.83) -- (2.00,-1.37);
\draw[transmitida] (0.25,-0.30) -- (2.13,-1.90);

\fill[black] (0,0) circle (2pt);
\draw[gray!80!black, line width=0.8pt] (0,0.62)
  arc[start angle=90,end angle=130,radius=0.62];
\node[fill=white, inner sep=1pt] at (-0.48,0.72) {$\theta_i$};
\draw[gray!80!black, line width=0.8pt] (0,0.62)
  arc[start angle=90,end angle=50,radius=0.62];
\node[fill=white, inner sep=1pt] at (0.55,0.72) {$\theta_r$};
\draw[gray!80!black, line width=0.8pt] (0,-0.62)
  arc[start angle=-90,end angle=-50,radius=0.62];
\node[fill=white, inner sep=1pt] at (0.47,-0.68) {$\theta_t$};

\node[align=center, text=blue!75!black, fill=white, inner sep=2pt]
  at (-3.2,1.35) {incidente\\P o S};
\node[align=center, text=teal!70!black, fill=white, inner sep=2pt]
  at (3.18,1.42) {reflejada\\P o S};
\node[align=center, text=orange!85!black, fill=white, inner sep=2pt]
  at (3.15,-1.55) {transmitida\\P o S};
\node[font=\footnotesize, text=gray!70!black, fill=white, inner sep=2pt]
  at (-2.92,-0.72) {frentes de onda};
\draw[gray!70!black, thin] (-2.16,-0.58) -- (-1.27,0.68);

\end{tikzpicture}
\par\medskip
\textbf{Figura 31.} Reflexión y transmisión de ondas en una interfaz entre dos medios. Las flechas indican la dirección de propagación y las líneas finas, los frentes de onda; los ángulos se miden respecto de la normal. Las ondas reflejadas y transmitidas pueden ser P o S, por lo que el esquema representa una posibilidad y no todas las ramas de conversión.
\end{figure}

\subsubsection{Ondas Rayleigh}

Las ondas Rayleigh son ondas superficiales que se propagan junto a una superficie libre. En un medio elástico homogéneo, las partículas cercanas a la superficie se mueven en el plano vertical que contiene la dirección de propagación y describen, de manera aproximada, órbitas elípticas retrógradas: el sentido del movimiento orbital es opuesto al avance de la onda.

\begin{figure}[H]
\centering
\begin{tikzpicture}[
  >=Latex,
  x=1.15cm,
  y=1.15cm,
  line cap=round,
  line join=round,
  font=\small,
  propagacion/.style={->, blue!75!black, line width=1.5pt},
  orbita/.style={draw=orange!85!black, line width=1.4pt}
]

% Medio elástico, superficie media y perfil de onda.
\fill[gray!8] (-5.8,-2.15) rectangle (5.8,0);
\draw[densely dashed, gray!55, line width=0.8pt] (-5.8,0) -- (5.8,0);
\draw[blue!55!black, line width=1.1pt]
  plot[smooth] coordinates {
    (-5.8,0) (-5.4,0.12) (-5.0,0.18) (-4.6,0.12) (-4.2,0)
    (-3.8,-0.12) (-3.4,-0.18) (-3.0,-0.12) (-2.6,0)
    (-2.2,0.12) (-1.8,0.18) (-1.4,0.12) (-1.0,0)
    (-0.6,-0.12) (-0.2,-0.18) (0.2,-0.12) (0.6,0)
    (1.0,0.12) (1.4,0.18) (1.8,0.12) (2.2,0)
    (2.6,-0.12) (3.0,-0.18) (3.4,-0.12) (3.8,0)
    (4.2,0.12) (4.6,0.18) (5.0,0.12) (5.4,0) (5.8,-0.12)
  };
\node[above left, text=gray!65!black, fill=white, inner sep=1pt]
  at (-5.55,0.06) {superficie media};

% Avance de la onda y órbita de una partícula superficial.
\draw[propagacion] (-4.8,0.95) -- (-2.3,0.95)
  node[midway, above, text=blue!70!black] {propagación};
\coordinate (P) at (1.15,0);
\draw[orbita, densely dashed] (P) ellipse (1.02 and 0.52);
\draw[->, orange!85!black, line width=1.5pt]
  (2.02,0.34) .. controls (1.82,0.52) and (1.53,0.52) .. (1.31,0.36);
\fill[orange!85!black] (1.15,0.52) circle (2.5pt);
\node[align=left, fill=white, inner sep=2pt] at (3.35,-0.78)
  {trayectoria elíptica\\retrógrada};
\draw[orange!85!black, thin] (2.65,-0.48) -- (2.10,-0.22);

% Ejes del plano vertical de propagación.
\draw[->, gray!75] (-4.8,-1.72) -- (-3.55,-1.72) node[right] {$x$};
\draw[->, gray!75] (-4.8,-1.72) -- (-4.8,-0.52) node[above] {$z$};

\end{tikzpicture}
\par\medskip
\textbf{Figura 29.} Esquema de una onda Rayleigh. La onda se propaga a lo largo de la superficie, mientras que una partícula superficial describe una órbita elíptica retrógrada en el plano vertical de propagación. El movimiento está ampliado.
\end{figure}

\subsection{Vibraciones radiales de una esfera elástica}

Consideremos una esfera elástica maciza de radio $R$. Por simetría esférica, un modo radial tiene desplazamiento
\[
\vec u(\vec r,t)=u_r(r,t)\,\hat r.
\]
La Figura 30 muestra la geometría y la dirección radial del movimiento.

\begin{figure}[H]
\centering
\begin{tikzpicture}[
  >=Latex,
  x=1.45cm,
  y=1.45cm,
  line cap=round,
  line join=round,
  font=\small,
  radial/.style={->, orange!80!black, line width=1.25pt},
  eje/.style={->, gray!70!black, line width=0.9pt}
]

\coordinate (C) at (0,0);
\fill[blue!5] (C) circle (2.0);
\draw[blue!65!black, line width=1.3pt] (C) circle (2.0);

\draw[eje] (C) -- (2.55,0) node[right] {$x$};
\draw[eje] (C) -- (0,2.55) node[above] {$z$};
\draw[eje] (C) -- (-1.25,1.05) node[above left] {$y$};

\draw[red!75!black, <->, line width=1.5pt] (C) -- (1.42,1.42)
  node[pos=0.72, above left, fill=white, inner sep=1.5pt] {$R$};
\foreach \angulo in {30,90,150,210,270,330} {
  \draw[radial]
    ({1.52*cos(\angulo)},{1.52*sin(\angulo)}) --
    ({2.28*cos(\angulo)},{2.28*sin(\angulo)});
}
\fill[black] (C) circle (1.8pt);

\end{tikzpicture}
\par\medskip
\textbf{Figura 30.} Corte esquemático de una esfera elástica de radio $R$. Las flechas indican el desplazamiento de un modo puramente radial.
\end{figure}

Para un sólido isotrópico y homogéneo, la ecuación de Navier es
\[
\rho\ddot{\vec u}
=\mu\nabla^2\vec u+(\lambda+\mu)\nabla(\nabla\cdot\vec u),
\]
donde $\lambda$ y $\mu$ son las constantes de Lamé. Un campo radial que depende solo de $r$ no tiene rotacional; por ello, podemos escribir
\[
\nabla\times\vec u=0,
\qquad
\vec u=\nabla\phi
=\frac{\partial\phi}{\partial r}\,\hat r.
\]
Al sustituir este potencial en la ecuación de Navier, se obtiene la ecuación de onda longitudinal
\[
c_L^2\nabla^2\phi=\ddot\phi,
\qquad
c_L^2=\frac{\lambda+2\mu}{\rho},
\qquad
c_T^2=\frac{\mu}{\rho}.
\]
Para oscilaciones armónicas proporcionales a $e^{-i\omega t}$, la ecuación radial queda
\[
\nabla^2\phi
=\frac{1}{r^2}\frac{d}{dr}\left(r^2\frac{d\phi}{dr}\right)
=-k^2\phi,
\qquad k=\frac{\omega}{c_L}.
\]
La ecuación radial es la ecuación de Bessel esférica de orden cero. Su solución general puede escribirse como
\[
\phi(r)=A j_0(kr)+B n_0(kr),
\qquad
j_0(x)=\frac{\sin x}{x},
\qquad
n_0(x)=-\frac{\cos x}{x}.
\]
La función $n_0$ diverge en el origen, de modo que la regularidad exige $B=0$. Absorbiendo el factor constante $1/k$ en la amplitud, queda
\[
\phi(r)=A\frac{\sin(kr)}{kr}.
\]

La ley de Hooke para un sólido isotrópico es
\[
\sigma_{ij}=2\mu U_{ij}+\lambda\delta_{ij}U_{kk},
\]
donde $U_{ij}$ es el tensor de deformación. Para el desplazamiento radial, $U_{rr}=\partial u_r/\partial r=\partial^2\phi/\partial r^2$ y $U_{kk}=\nabla\cdot\vec u=\nabla^2\phi$. Así,
\[
\sigma_{rr}
=2\mu\frac{d^2\phi}{dr^2}+\lambda\nabla^2\phi
=\rho\left\{(c_L^2-2c_T^2)\nabla^2\phi
+2c_T^2\frac{d^2\phi}{dr^2}\right\}.
\]
Usando la ecuación de Helmholtz, se obtiene
\[
\frac{\sigma_{rr}}{\rho}
=-\omega^2\phi-\frac{4c_T^2}{r}\frac{d\phi}{dr}.
\]

En una superficie libre, el esfuerzo radial se anula: $\sigma_{rr}(R)=0$. Definamos $q=kR$. Para la solución radial,
\[
\phi(R)=A\frac{\sin q}{q},
\qquad
\phi'(R)=\frac{A}{Rq}(q\cos q-\sin q).
\]
Al sustituir en la condición de frontera, resulta
\[
\omega^2R^2\sin q+4c_T^2(q\cos q-\sin q)=0.
\]
Reordenando y usando $q=kR$ y $\omega=kc_L$, la ecuación trascendental para las frecuencias características queda
\[
\frac{\tan(kR)}{kR}
=\frac{1}{1-\left(\dfrac{kR c_L}{2c_T}\right)^2}.
\]
Sus raíces determinan los modos normales radiales de la esfera.

\textbf{Conexión con las oscilaciones radiales estelares.}
En ambos problemas se estudian pequeñas perturbaciones alrededor de un equilibrio esférico y se obtienen modos radiales con frecuencias características. Para la esfera elástica, la condición de frontera es $\sigma_{rr}(R)=0$; para una estrella fluida, la condición análoga es que la perturbación lagrangiana de presión se anule en la superficie, $\Delta p(R)=0$. Un sólido elástico soporta ondas longitudinales y transversales, mientras que un fluido ideal no soporta esfuerzos cortantes ($\mu=0$) y solo admite perturbaciones compresionales.

\subsection{Vibraciones radiales de una cavidad esférica en un medio elástico infinito}

Ahora consideremos una cavidad de radio $R$ rodeada por un medio elástico que se extiende indefinidamente. Cuando la cavidad oscila, emite ondas longitudinales hacia el exterior; estas transportan energía y amortiguan la oscilación. En el régimen $c_L\gg c_T$, la radiación es débil y el amortiguamiento es pequeño.

Como $K=\lambda+2\mu/3$, $c_L^2=(K+4\mu/3)/\rho$ y $c_T^2=\mu/\rho$, este régimen corresponde a un módulo de compresibilidad $K$ mucho mayor que el módulo de cizalladura $\mu$. La oscilación considerada requiere $\mu\ne0$; si el medio no ofrece resistencia al cizallamiento, desaparece este modo elástico de cavidad.

\begin{figure}[H]
\centering
\begin{tikzpicture}[
  >=Latex,
  x=1.25cm,
  y=1.25cm,
  line cap=round,
  line join=round,
  font=\small,
  frente/.style={draw=blue!65!black, line width=1.25pt},
  propagacion/.style={->, orange!85!black, line width=1.35pt},
  dimension/.style={<->, black!75, line width=0.9pt}
]

\coordinate (C) at (0,0);

% Frentes de onda en el corte radial del medio.
\foreach \radio in {1.35,1.95,2.55,3.15} {
  \draw[frente] (C) circle (\radio);
}

% Cavidad y radios característicos.
\fill[blue!3] (C) circle (0.86);
\draw[black!70, line width=1.4pt] (C) circle (0.86);
\node[font=\small\bfseries] at (0,0.28) {cavidad};
\fill[black!70] (C) circle (1.5pt);
\draw[dimension] (C) -- (0.78,-0.36)
  node[pos=0.66, below=1pt, fill=white, inner sep=1.5pt] {$R$};
\draw[dimension] (C) -- (-2.12,-2.12)
  node[pos=0.72, below left, fill=white, inner sep=1.5pt] {$r$};

% Flechas radiales: propagación hacia afuera.
\foreach \angulo in {30,90,150,270,330} {
  \draw[propagacion]
    ({1.05*cos(\angulo)},{1.05*sin(\angulo)}) --
    ({1.62*cos(\angulo)},{1.62*sin(\angulo)});
}

% Identificación de las ondas y del material.
\node[align=center, text=blue!55!black, fill=white, inner sep=2pt]
  at (-4.35,2.10) {frentes de onda\\longitudinal saliente};
\draw[blue!55!black, thin] (-3.42,1.80) -- (-2.05,2.05);
\node[align=center, fill=white, inner sep=2pt] at (4.15,-2.05)
  {medio elástico\\infinito};

\end{tikzpicture}
\par\medskip
\textbf{Figura 32.} Corte esquemático de una cavidad esférica en un medio elástico infinito. Los círculos representan frentes de onda longitudinal que se propagan radialmente hacia el exterior; $R$ es el radio de la cavidad y $r$ la distancia al centro.
\end{figure}

Buscamos una onda esférica saliente,
\[
\phi(r)=\frac{A e^{ikr}}{r},
\qquad k=\frac{\omega}{c_L}.
\]
Al imponer la condición de superficie libre, $\sigma_{rr}(R)=0$, resulta
\[
\left(\frac{kR c_L}{c_T}\right)^2=4(1-ikR).
\]
Para $c_L\gg c_T$, la frecuencia compleja se aproxima por
\[
\omega\simeq\frac{2c_T}{R}\left(1-i\frac{c_T}{c_L}\right).
\]
La parte real determina la frecuencia de oscilación y la parte imaginaria describe el amortiguamiento. En el límite $c_L\to\infty$, este amortiguamiento por radiación desaparece. En esta aproximación, $kR\simeq2c_T/c_L\ll1$, así que la longitud de onda es mucho mayor que el radio de la cavidad. En contraste, para una esfera maciza y $c_L\gg c_T$, la primera frecuencia característica corresponde a $kR\simeq\pi$.

\subsection{Ondas sísmicas en un medio elástico}

Las ondas sísmicas ilustran cómo se propagan las perturbaciones elásticas desde un foco. Las ondas P son longitudinales y las más rápidas; las ondas S son transversales y no se propagan en fluidos. Las ondas superficiales viajan cerca de la superficie terrestre y pueden producir movimientos intensos. Por ello, la detección de ondas P puede dar aviso antes de la llegada de ondas posteriores, aunque el tiempo disponible depende de la distancia y del sistema de alerta.

\begin{figure}[H]
\centering
\begin{tikzpicture}[
  >=Latex,
  x=1.25cm,
  y=1.2cm,
  line cap=round,
  line join=round,
  font=\small,
  ondaP/.style={->, blue!75!black, line width=1.6pt},
  ondaS/.style={->, orange!85!black, line width=1.6pt},
  ondaSup/.style={->, violet!75!black, line width=1.6pt}
]

% Corte esquemático del terreno.
\fill[gray!8] (-5.8,-1.85) rectangle (5.8,0.8);
\draw[black!75, line width=1.3pt] (-5.8,0.8) -- (5.8,0.8);
\node[above left, text=gray!65!black] at (-5.6,0.8) {superficie};

% Foco, epicentro y estación receptora.
\coordinate (foco) at (0,-1.3);
\coordinate (epicentro) at (0,0.8);
\coordinate (estacion) at (4.6,0.8);
\draw[densely dashed, gray!65] (foco) -- (epicentro);
\fill[red!75!black] (foco) circle (3pt);
\node[below left, align=right, text=red!70!black] at (-0.08,-1.3)
  {foco sísmico};
\fill[red!75!black] (epicentro) circle (2pt);
\node[above, text=red!70!black] at (0,0.88) {epicentro};

% Ondas internas P y S, y onda superficial.
\draw[ondaP] (foco) -- (estacion);
\draw[ondaS] (foco) .. controls (1.50,-1.10) and (3.10,-0.15) .. (estacion);
\draw[ondaSup] (0.18,0.81)
  .. controls (0.62,1.24) and (1.05,0.38) .. (1.48,0.81)
  .. controls (1.90,1.24) and (2.33,0.38) .. (2.76,0.81)
  .. controls (3.18,1.24) and (3.60,0.38) .. (4.18,0.81);

% Instrumento situado en la superficie.
\draw[black!70, line width=1pt] (4.6,0.8) -- (4.6,1.12);
\draw[black!70, line width=1pt] (4.42,1.12) -- (4.78,1.12);
\node[above right, align=left] at (4.82,1.38) {estación\\sísmica};

% Leyenda alineada bajo el corte.
\draw[ondaP] (-5.1,-2.3) -- (-4.35,-2.3);
\node[anchor=west, font=\footnotesize] at (-4.2,-2.3) {P: compresional};
\draw[ondaS] (-1.15,-2.3) -- (-0.4,-2.3);
\node[anchor=west, font=\footnotesize] at (-0.25,-2.3) {S: transversal};
\draw[ondaSup] (2.55,-2.3) -- (3.3,-2.3);
\node[anchor=west, font=\footnotesize] at (3.45,-2.3) {superficial};

\end{tikzpicture}
\par\medskip
\textbf{Figura 33.} Esquema de las trayectorias de ondas sísmicas desde un foco hasta la superficie. Las ondas P atraviesan el interior como perturbaciones compresionales; las ondas S lo hacen como perturbaciones transversales, mientras que las ondas superficiales se propagan junto al límite del terreno.
\end{figure}